\chapter{Conclusion}
\label{ch:Conclusion}

Consistency preservation rules in the Reactions Language establish a fixed mapping between metamodels.
Enabling a second mapping, for example, representing UML classes as Java interfaces, necessitates maintaining a separate set of rules.
A persisted V-SUM does not retain information about which rules were used to derive its models.
When the rules are modified, the models retain the state established by the previous rules.
Both of these challenges are addressed in this thesis.
Feature annotations place the reactions of all mappings in one 150\% model.
A preprocessor derives the reactions of a single configuration.
A migration leads a persisted V-SUM to the changed rules.
A full migration re-derives all other models from a dominant model, whereas a selective migration only repropagates the affected elements.
A preservation step reports or restores the underived content that both leave behind.
Overall, the feature annotations achieve their intended purpose, although some drawbacks remain.
The migration process also achieves its objective, and selective migration confines the required work to a limited portion of the V-SUM.
Some manual work remains, however, and the selective migration is hardly faster than a full one.

\section{Feature annotations}
\label{sec:Conclusion:FeatureAnnotations}

The process of annotating the umljava reactions yielded notable insights.
The 150\% model has 3,266 lines of code.
Maintaining the nine configurations as separate copies would require 23,791 lines.
It is evident that even two separate copies would exceed the dimensions of the 150\% model.
A comparison of the 150\% model with the original rules reveals a 36.4\% growth, despite the former encompassing nine configurations, in contrast to the one configuration covered by the latter.
The growth consists primarily of one annotation per reaction and the reactions and routines of the new features.
Modifications to features and configurations are inexpensive as well.
A configuration is defined as a single JSON file, while a feature encompasses one to five reaction files.
The derived rules also demonstrate a similar behavior to the original ones.
Config1 passes in all 162 of the original tests from the case study.
In all nine configurations, the 189 tests exhibit the anticipated behavior with the feature gates.
The preprocessing stage requires minimal additional time.
The duration of this process is 0.628 seconds, which constitutes 0.55\% of the total time allocated for preprocessing, building, and testing a configuration.
The preprocessor contains no code for a particular rule set or metamodel.
Consequently, the transfer occurs to other rules, and all 52 executable tests of pcmumlclass continue to pass after annotating and preprocessing.

However, it is important to note that this approach is not without its drawbacks.
It has been demonstrated that features have the capacity to interact through shared routines.
It was only after migrating every pair of configurations that several such interactions became apparent.
Furthermore, it is important to note that the preprocessor does not verify a configuration against the feature model.
The responsibility for ensuring the validity of these methods therefore falls upon the methodologist.
In addition, the pcmumlclass transfer utilizes a single feature.
This finding indicates that the mechanism functions under alternative rules, yet it does not demonstrate that these alternative rules are enhanced by variability.

\section{Migration}
\label{sec:Conclusion:Migration}

Both types of migration successfully managed each configuration switch in the library example.
All 72 selective migrations were completed without reverting to a full migration.
Full migrations were completed for all 81 pairs using each strategy.

A full migration may determine its dominant model using three distinct approaches.
The explicit strategy has minimal costs and ensures predictable outcomes.
The reachability strategy is similarly cost-effective, but it does not independently determine the dominant model.
In the library example, UML and Java are mutually reachable, resulting in ties across all 81 configuration pairs.
In the BrakeCaseStudy, a tie occurred between CAD and Simulink, with CAD selected as the model listed first.
The rules applied in that case study do not enable re-derivation of other models from CAD.
Given the presence of only two metamodels, the umljava rules are inadequate for evaluating this strategy.
The fewest changes strategy adjusts the selection based on each configuration pair.
In every pair within the library example, this strategy selected the less costly option.
This approach reduced the median number of changes from 264 for UML to 185.
In the BrakeCaseStudy, the unsuccessful trial of this strategy excluded CAD.
These trials increase computational time, consuming 61\% of a migration duration, resulting in 8.9 seconds instead of 5.4 seconds.
Regardless of the chosen strategy, a full migration re-derives each derived model, even if the reactions remain unchanged.

Selective migration prevents this issue.
The process of identifying the impact of a configuration switch is effective.
Comparing identifiers is sufficient because no rule shared between two configurations has altered its semantic hash.
Among the locations where a switch introduces inconsistencies, 96.1\% are contained within affected elements.
The resulting scope remains limited, with a median of 8.1\% of the rules and 1.42\% of the model elements affected.
For unchanged rules, selective migration completes in 49 milliseconds, while a full migration re-derives all elements in 5.1 seconds.
Selective migration also preserves content that remains unchanged.
For example, in 22 pairs, the handwritten body of \texttt{Member.totalWeight} is retained, whereas a full migration would discard it.

However, reducing the scope does not significantly decrease the runtime.
Selective migration requires a median time of 5.1 seconds, compared to 5.4 seconds for a full migration.
In 16 out of 72 pairs, selective migration is slower, by at most 9\%.
For a V-SUM of this scale, the majority of the runtime is consumed by model loading and the preservation step.
Both processes require approximately the same amount of time in either migration scenario.
Additionally, the trigger summaries omit \texttt{check} and \texttt{with} guards, resulting in unnecessary classifier repropagation in 48 out of 72 migrations.
Furthermore, selective migration necessitates more supporting infrastructure than a full migration.
It needs a registry that every change propagation persists, as well as semantic hashes and trigger summaries from the compiler.
The evaluation does not utilize the semantic hashes, as configuration switches do not alter any shared rules.

Switches that do not deselect any of the five features contributing to the UML model result in a state equivalent to a derivation from scratch.
Switching from config1, config7, or config8 cannot deselect any features, since these configurations select none of them.
Similarly, switching to config9, which selects all five features, does not result in any deselection.
No compilation errors arise from elements that are incorrectly derived by the new rules.
A limitation persists for features that also modify the dominant UML model.
Examples are the prefix of \texttt{IBorrowable} and private attributes.
Repropagation reinserts the dominant elements unchanged, and no rule in the new configuration reverses this effect.
A full migration also retains the dominant model and is subject to the same limitation.
Consequently, 29 out of the 72 selective migrations are equivalent to a derivation from scratch.
Round trips to another configuration and subsequent return restore the original V-SUM in only 2 out of 72 cases.
Thirty of the 72 migrations compile successfully, while the remaining cases fail due to the absence of the handwritten body for \texttt{Member.totalWeight}.

The preservation step identifies handwritten content in all instances where a migration process removes such material.
The report documented the complete presence of the body of \texttt{Member.totalWeight} in 27 migrations and its partial presence in seven migrations.
In the eight migrations into config7, it marks the body as lost, since an interface method holds no statements.
The reports also enumerate 957 items that were not authored by individuals, such as modifiers and initial values lacking direct correspondence.
Restoring these items using \texttt{--preserve user} would place previous values alongside the newly generated ones.
For the library example, \texttt{--preserve report} represents the appropriate level of preservation, requiring manual reintegration of the handwritten body.

The migration process operates largely independently of specific metamodels.
Ten metamodels using XMI need no code of their own.
Only the Java implementation requires an adapter, consisting of 190 lines of code supporting eleven hooks.
JaMoPP stores models as source text, resolves references using a global classpath, and does not assign identifiers to elements.
Consequently, the umljava rules presented a demanding example because the requirements would not have been apparent from XMI metamodels alone.

In summary, migration enables a persisted V-SUM to function effectively under changed rules.
However, certain manual tasks are still required.
Handwritten content must be restored from the preservation report manually.
Residual effects of deselected features in the dominant model must be eliminated.
